\chapter{Axiomes de la topologie}
\lecture{28 septembre}

\section{Topologie et base}

\begin{definition}[Espace topologique] \label{def: Espace topologique}
	Soit $X$ un ensemble, une \textbf{topologie} sur $X$ est un sous-ensemble $\T \subset \mathcal P(X)$ tel que
	\begin{enumerate}
		\item[(T1)] \label{def: T1} $\empty, X \in \T$
		\item[(T2)] \label{def: T2} $\forall \mathcal A \subset \T, \quad \bigcup_{U \in \mathcal A} U \in \T$
		\item[(T3)] \label{def: T3} $\forall U_1, \ldots, U_n \in \T, \quad \bigcup_{i=1}^n U_i \in \T$
	\end{enumerate}
	le couple $(X, \T)$ est alors appelé un \textbf{espace topologique} et les éléments de $\T$ sont appelés des \textbf{ouverts} de $(X, \T)$.
\end{definition}

\begin{exemple}
	Tout ensemble $X$ possède les topologies suivantes
	\begin{enumerate}
		\item La \textbf{topologie grossière} définit par
		\[
		\Tgr = \{ \emptyset, X\}
		\]
		\item La \textbf{topologie discrète} définit par
		\[
		\Tdis = \mathcal P(X)
		\]
		\
	
	 \item Mais il existe un topologie moins intuitive qu'on appelle \textbf{topologie du complément fini} que l'on définit par
	\[
	\Tcomp = \{ U \in \mathcal P(X) \mid \# U^c < \infty\} \cup \{\emptyset\}
	\]
	En effet on a 
	\begin{enumerate}
		\item[\TU : ] Par définition $\emptyset \in \Tcomp$. De plus $X \in \Tcomp$ car $\# X^c = \# \emptyset = 0 < \infty$
		
		\item[\TD : ] De plus pour $ \mathcal A \subset \Tcomp$ alors
		\[
		\bigg(\bigcup_{U \in \mathcal A} U \bigg)^c = \bigcap_{U \in \mathcal A} U^c
		\]
		donc
		\[
		\bigg(\bigcup_{U \in \mathcal A} U \bigg)^c = \#\bigcap_{U \in \mathcal A} U^c \leqslant \# V^c < \infty
		\]
		où $V \in \mathcal A$ quelconque. 
		
		\item[\TT: ] Et finalement pour $U_1, \ldots, U_n \in \Tcomp$ alors
		\[
		\bigg(\bigcap_{i=1}^n U_i \bigg)^c = \bigcup_{i=1}^n U_i^c
		\]
		ainsi
		\[
		\#\bigg(\bigcap_{i=1}^n U_i \bigg)^c = \#\bigcup_{i=1}^n U_i^c \leqslant \sum_{i=1}^n \# U_i^c < \infty
		\]
		Ainsi $(X, \Tcomp)$ est bien une topologie.
		
	\end{enumerate}
\end{enumerate}

\end{exemple}

\begin{definition}
	Soit $\T, \T'$ des topologies sir $X$. Si $\T \subset \T'$, alors on dit que $\T'$ est plus \textbf{fine} que $\T$. Mais aucune des topologies n'est incluses dans l'autre alors on dira qu'elles sont \textbf{incomparable}.
\end{definition}

\subsection{Bases}

\begin{definition}[Base topologique] \label{def: Base topologique}
	Soit $X$ un ensemble. Une \textbf{base de topologie} est un sous-ensemble $\mathcal B \subset \mathcal P(X)$ qui vérifie
	\begin{enumerate}
		\item[(B1)] \label{def: B1} On a $X = \bigcup_{B \in \mathcal B} B$,
		\item[(B2)] \label{def: B2} Pour tout $B_1, B_2 \in \mathcal B$, on a que pour tout $x \in B_1 \cap B_2$, il existe un $B_3 \in \mathcal B$ tel que 
		\[
		x \in B_3 \subset B_1 \cap B_2
		\]
		Il est facile de voir qu'une topologie $\T$ sur $X$ est aussi une base de topologie. Par contre la réciproque est généralement fausse.
	\end{enumerate}
\end{definition}

\begin{figure}[H]
	\centering
	\includegraphics[width=0.75\textwidth]{"images/base topologique.png"}
	\caption{Illustration honteusement volé dans le polycopié du cours}
	\label{fig: Illustration d'une base topologique}
\end{figure}

\begin{definition} \label{def: Topologie engendrée}
	Soit $\B$ une base de topologie sur $X$, alors on définit la \textbf{topologie engendrée} par $\B$ comme
	\[
	\T_\B = \left\{ \bigcup_{\mathcal V \in \mathcal A} V \mid A \subset \B \right\}
	\]
\end{definition}

\begin{lemme}
	$\T_\B$ est bien une topologie sur $X$.
\end{lemme}

\prf{On doit vérifié que $\TB$ vérifie les trois conditions requises pour être une topologie.
	\begin{enumerate}
		\item[\TU :] Par définition d'une base on a en particulier pour $\mathcal A = \mathcal B$
		\[
		X = \bigcup_{V \in \mathcal B} V \in \T_\B
		\]
		et de même pour $\mathcal A = \emptyset$ on a bien $\emptyset \in \T_\B$.
		
		\item[\TD :] Soit $\{U_\alpha \mid \alpha \in I\} \subset \T_\B$, alors par définition de $\T_\B$ pour tout $\alpha \in I$ il existe $\mathcal A_\alpha \subset \B$ tel que $U_\alpha = \bigcup_{V \in \mathcal A_\alpha} V$, d'où
		\[
		\bigcup_{\alpha \in I} U_\alpha = \bigcup_{\alpha \in I} \bigcup_{V \in \mathcal A_\alpha} V = \bigcup_{V \subset \bigcup_{\alpha \in I} \mathcal A_\alpha} V \in \T_\B
		\]
		
		\item[\TT :] Finalement, soit $\{U_1, \ldots, U_n\} \subset \T_\B$, alors par définition de $\T_\B$ pour tout $i$ il existe $\mathcal A_i \subset \B$ tel que $U_i = \bigcup_{V \in \mathcal A_i} V$ ainsi
		\[
		\bigcap_{i=1}^n U_i = \bigcap_{i=1}^n \bigcup_{V \in \mathcal A_i} V
		\]
		Pour conclure, il suffit de considérer le cas $n = 2$ car on peut ensuite conclure par récurrence, car si on sait que
		\[
		U_1, U_2 \in \T_\B \implies U_1 \cap U_2 \in \T_\B
		\]
		alors
		\[
		U_1 \cap U_2 \cap U_3 = (U_1 \cap U_2) \cap U_3 \in \T_\B
		\]
		Or
		\[
		U_1 \cap U_2 = \bigcup_{V \in A_1} V \cap \bigcup_{V' \in A_2} V' = \bigcup_{V \in \mathcal A_1, \, V \in \mathcal A_2} (V \cap V')
		\]
		or par définition de la base on a
		\[
		V \cap V' = \bigcap_{x \in V \cap V'} V_x
		\]
		d'où 
		\[
		U_1 \cap U_2 = \bigcup_{V \in \mathcal A_1, \, V \in \mathcal A_2}\bigcap_{x \in V \cap V'} V_x \in \T_\B
		\]
	\end{enumerate}
	Ainsi la topologie engendrée $\TB$ est bien une topologie.
}

\begin{theoreme}[Caractérisation de la topologie engendrée]
	Soit $\B$ une base de topologie sur un ensemble $X$, alors
	\[
	\T_\B = \{U \in \mathcal P(X) \mid \forall x \in U, \, \exists B \in \B, \, x \in B \subset U\}
	\]
\end{theoreme}

\prf{Soit $U = \bigcup_{V\in \mathcal A} V \in \T_\B$, alors pour tout $x \in U$ il existe $v \in \mathcal A \subset B$ tel que $x \in V$. Ainsi $\T_\B \subset \T'$.
	
	Réciproquement, soit $U \in \T'$, alors pour tout $x \in U$, il existe $V_x$ tel que $x \in V_x \subset U$ et donc par conséquent $U = \bigcup_{x \in U} V_x \in \T_\B$.
}

\lecture{5 octobre}

\begin{lemme}
	Soit $(X,d)$ un espace métrique, on définit les boules ouvertes
	\[
	B_d(x,r) = \{y \in X \mid d(x,y) < r\}
	\]
	pour tout $x \in X$ et $r > 0$. Alors
	\[
	\mathcal B_d = \{ B_d(x,r) \mid x \in X, ~ r > 0\}
	\]
	est une base de topologie sur $X$.
\end{lemme}

\prf{
\begin{enumerate}
	\item[\BU :] On sait que pour tout $x \in X$, on a $x \in B_d(x,r)$ pour $r > 0$, ce qui implique que
	\[
	X = \bigcup_{x \in X} B_d(x,r)
	\]
	\item[\BD :] De plus soit $B_d(x_1, r_1)$ et $B_d(x_2, r_2)$ deux boules, alors si on considère
	\[
	r = \min\{ r_1 - d(x, x_1), r_2 - d(x, x_2)\}
	\]
	alors $B_d(x,r) \subset B_d(x_1, r_1) \cap B_d(x_2, r_2)$ car si $y \in B_d(x,r)$ alors $d(x,y) < r$ et donc
	\[
	d(x_i, y) \leqslant d(x_1, x) + d(x,y)  < d(x_i, x) + (r_i - d(x_i, x)) < r
	\]
	pour $i =1,2$. 
\end{enumerate}
Ainsi on pose $\mathcal T_d = \mathcal T_{\mathcal B_d}$ la \textbf{topologie métrique} sur $X$ induite par la métrique $d$. Par la suite si $X = \Rn$ et $d = d_{euc}$, on écrit $\T_{st} = \T_{d_{euc}}$ la \textbf{topologie standard}.
}

\begin{lemme}[Comparaison de bases] \label{lem: Comparaison de bases}
	Soit $X$ un ensemble, soient $\B_1, \B_2$ des bases topologies sur $X$. Alors
	\[
	\T_{\B_1} \subset \T_{\B_2} \iff \forall B_1 \in \B_1, \, \forall x \in B_1, \, \exists B_2 \in \mathcal B_2, \quad x \in B_2 \subset B_1
	\]
\end{lemme}

\prf{
	\begin{enumerate}
		\item[$\implies$ :] Puisque $\mathcal B_1 \subset \T_{\B_1} \subset \T_{\B_2}$, alors pour tout $B_1 \in \B_1$, on a que pour tout $x \in B_1$, il existe $B_2 \in \B_2$ tel que $x \in B_2 \subset B_1$.
		
		\item[$\impliedby$ :] Réciproquement, alors par hypothèse on a $B_1 \in \T_{\B_2}$ pour tout $B_1 \in \B_1$. Autrement dit, $\B_1 \subset \T_{\B_2}$. Soit $U \in \T_{\B_1}$, alors par la \autoref{def: Topologie engendrée}, il existe une collection $\{B_i\} \subset \B_1$ telle que $U = \bigcup_{i \in I} B_i$. Or on sait que $B_i \in \T_{\B_2}$ pour tout $i \in I$. Puisque $\T_{B_2}$ est une topologie, alors
		\[
		U = \bigcup_{i \in I} B_i \in \T_{\B_2}
		\]
	\end{enumerate}
}

\begin{lemme}[Test de base] \label{lem: Test de base}
Soit $(X, \T)$ un espace topologique, soit $\mathcal C \subset \T$, alors $\mathcal C$ est une base de topologie et $\T_{\mathcal C} = \T$ si et seulement si
\[
\forall U \in \T, \, \forall x \in U, \, \exists C \in \mathcal C, \quad x \in C \subset U
\]
\end{lemme}

\prf{
	\begin{enumerate}
		\item[$\implies$ :] Soit $\mathcal C$ une base de topologie, et $\T_{\mathcal C} = \T = \T_\T$, et donc en particulier $\T_\T \subset \T_\C$, d'où par le \autoref{lem: Comparaison de bases},
		\[
		\forall U \in \T, \, \forall x \in U, \, \exists C \in \C, \quad x \in C \subset U
		\]
		
		\item[$\impliedby$ :] On doit montrer que $\C$ engendre $\T$, il faut donc vérifier les deux axiomes
		\begin{enumerate}
			\item[\BU :] par hypothèse sur $\C$, en particulier pour $X \in \T$, 
			\[
			\forall x \in X, \, \exists C \in \C, \quad x \in C \subset X
			\]
			d'où le recouvrement.
			\item[\BD :] Soit $C_1, C_2 \in \C$ et soit $x \in C_1 \cap C_2$, mais comme $C_1$ et $C_2$ sont des ouverts, alors $C_1 \cap C_2$ est aussi un ouvert, et donc par hypothèse on a pour $x \in C_1 \cap C_2$, et on en déduit qu'il existe $C_3 \in \C$ tel que
			\[
			x \in C_3 \subset C_1 \cap C_2
			\]
			Il nous à montrer que $\T \subset \T_\C$. Les hypothèses sur $\C$ implique, par le \autoref{lem: Comparaison de bases} que $\T = \T_\T \subset \T_\C$.
		\end{enumerate}
	\end{enumerate}
}

\section{Les sous-bases de la topologie}

\begin{definition}
	Soit $X$ un ensemble. Une \textbf{sous-base de topologie} $\SB$ est un sous-ensemble de $\mathcal P(X)$ tel que
	\begin{enumerate}
		\item[(S1)] \label{def: S1} $X = \bigcup_{S \in \SB} S$
	\end{enumerate}	
\end{definition}

\begin{lemme}
	Étant donné une sous-base de topologie sur $X$, alors la collection
	\[
	\B_\SB = \{ S_1 \cap \ldots  \cap S_n \mid S_i \in \SB, \, n \geqslant 1\}
	\]
	est une base de topologie. 
\end{lemme}

\prf{On doit montrer que $\B_\SB$ vérifie bien les axiomes d'une base
	\begin{enumerate}
		\item[\BU :] Comme $\SB$ est une sous-base, elle vérifie \SU = \BU. Comme $\SB \subset \BS$, on peut conclure que $\BS$ vérigie \BU.
		
		\item[\BD :] Soient $B = \bigcap_{i=1}^m S_i$ et $\tilde B = \bigcap_{j=1}^n \tilde S_j$ deux éléments de $\BS$. Alors leur intersection
		\[
		B \cap \tilde B = \left(\bigcap_{i=1}^m S_i\right) \cap \left(\bigcap_{j=1}^n \tilde S_j\right) \in \BS
		\]
	\end{enumerate}
	Ainsi $\BS$ est bien une base.
}	
\begin{figure}[H]
	\centering
	\includegraphics[width=0.5\textwidth]{"images/diagramme topologique.png"}
	\caption{Illustration honteusement volé dans le polycopié du cours}
	\label{fig: Diagramme topologique}
\end{figure}


\section{Autour de la notion de fermé}

\begin{definition}[Ensemble fermé] \label{def: Ensemble fermé}
	Soit $(X, \T)$ un espace topologique. Soit $A \subset X$, alors $A$ est dit \textbf{fermé par rapport à la topologie} $\T$ si $A^c \in \T$.
\end{definition}

\begin{exemple}
	On a
	\begin{enumerate}
	 \item Par si on considère la topologie grossière $(X, \T_{gr})$ alors les seuls fermés sont $X$ et $\emptyset$ donc $\emptyset$ et $X$ sont à la fois ouvert et fermé. 
	
	Mais plus généralement, pour tout espace topologique $\T$ sur $X$ alors $\emptyset$ et $X$ sont ouverts et fermés.
	
	 \item Maintenant si on considère la topologie discrète $\T_{disc}$ alors tout $A \subset X$ est fermé et ouvert.
	\end{enumerate}
\end{exemple}

\begin{proposition}
	Soit $(X, \T)$ un espace topologique, alors
	\begin{enumerate}
		\item $\emptyset, \, X$ sont fermés par rapport à $\T$
		\item Toute intersection de fermés est fermé.
		\item Toute réunion finie de fermés est fermé.
	\end{enumerate}
\end{proposition}

\prf{
Soit $\{A_\alpha \mid \alpha \in I\}$ une collection de fermés de $X$ par rapport à $\T$, alors
\[
\left(\bigcap_{\alpha \in I} A_\alpha\right)^c = \bigcup_{\alpha \in I} A_\alpha^c \in \T
\]
donc $\bigcap A_\alpha$ est bien fermé.

De plus, soient $A_1, \ldots, A_n$ des fermés par rapport à $\T$, alors
\[
\left(\bigcup_{i = 1}^n A_i\right)^c = \bigcap_{i=1}^n A_i^c \in \T
\]
Donc $\bigcup A_i$ est bien fermé
}

\begin{definition}
	Soit $(X, \T)$ un espace topologique, et $A \subset X$, alors on définit
	\begin{enumerate}
		\item L'\textbf{intérieur} de $A$ comme
		\[
		\Int A := \mathring A := \bigcup_{\{U \mid U \in \T, \, U \subset A\}} U
		\]
		\item L'\textbf{adhérence} de $A$ comme
		\[
		\overline A = \bigcap_{ \{C \mid C \text{ fermé}, \, A \subset C\}} C 
		\]
	\end{enumerate}
\end{definition}

\begin{remarque}
	On a toujours les inclusions
	\[
	\Int A \subset A \subset \overline A 
	\]
	avec des égalités si et seulement $A$ est respectivement ouvert ou fermé.
\end{remarque}


\begin{theoreme}
	Soient $(X, \T)$ un espace topologique et $\B$ une base de $\T$. Alors pour tout $A \subset X$,
	\[
	\overline A = \{x \in X \mid \forall B \in \B, \, x \in B, \quad A \cap B \neq \emptyset\}
	\]
\end{theoreme}

\prf{Soit $x \in \overline A$, alors pour tout fermé $C$ tel que $A \subset C$ alors $x \in C$. De plus soit $V \in \B$ tel que $x \in \B$, si on suppose que $A \cap B = \emptyset$, alors on aurait
	\[
	(A \cap B)^c = \emptyset^c = X
	\]
	or $(A \cap B)^c = A^c \cup B^c$. Mais comme $A \subset X$ et $A \cap A^c = \emptyset$, on en déduit que $A \subset B^c$ et donc on devrait avoir $x \in B^c$ ce qui contradictoire. Par conséquent, on a bien $A \cap B \neq \emptyset$.
}